% !TEX root = PSanté_03_présentation.tex

\usefonttheme{serif}
\usepackage{fontspec}
	\setmainfont{TeX Gyre Heros}
\usepackage{unicode-math}
\usepackage{lualatex-math}
	\setmathfont{TeX Gyre Termes Math}
\usepackage{polyglossia}
\setdefaultlanguage[frenchpart=false]{french}
\setotherlanguage{english}
\usepackage{microtype}
\usepackage[locale = FR,
            separate-uncertainty,
            multi-part-units = single,
            range-units = single]{siunitx}
	\DeclareSIUnit\an{an}
  \DeclareSIUnit{\octet}{o}
\usepackage{amsmath}
\usepackage{amsfonts}
\usepackage{amssymb}
\usepackage{array}
\usepackage{graphicx}
\graphicspath{{./Figures/}}
\usepackage{booktabs}
\usepackage{tabularx}
\usepackage{multirow}
\usepackage{multicol}
    \newcolumntype{L}{>{\raggedright\arraybackslash}X}
    \newcolumntype{R}{>{\raggedleft\arraybackslash}X}
\usepackage{tikz}
\usetikzlibrary{mindmap}
\usetikzlibrary{overlay-beamer-styles}
\usepackage{subcaption}
\usepackage[]{animate}
\usepackage{float}
\usepackage{csquotes}

\usetheme[secheader
         ]{Boadilla}
\usecolortheme{seagull}
\setbeamertemplate{enumerate items}[default]
\setbeamertemplate{itemize items}{-}
\setbeamertemplate{navigation symbols}{}
\setbeamertemplate{bibliography item}{}
\setbeamerfont{framesubtitle}{size=\large}
\setbeamertemplate{section in toc}[sections numbered]
%\setbeamertemplate{subsection in toc}[subsections numbered]

\title[Concevez une application au service de la santé publique]{Projet 3 : Concevez une application au service de la santé publique}
\author[Lancelot \textsc{Leclercq}]{Lancelot \textsc{Leclercq}} 
\institute[]{}
\date[]{\small{2 novembre 2021}}

\AtBeginSection[]{
  \begin{frame}
  \vfill
  \centering
    \usebeamerfont{title}\insertsectionhead\par%
  \vfill
  \end{frame}
}

\begin{document}
\setbeamercolor{background canvas}{bg=gray!20}
\begin{frame}[plain]
  \titlepage
\end{frame}

\section{Introduction}

\subsection{Problématique}
\begin{frame}{\insertsubsection}
  \begin{columns}
    \begin{column}{.6\textwidth}
      \begin{itemize}
        \item Application au service de la santé publique
        \item[]
        \item Idées d'application en lien avec l'alimentation
      \end{itemize}
    \end{column}
    \begin{column}{.4\textwidth}
      \includegraphics[width=.8\textwidth]{./Sante-publique-France-logo.pdf}
    \end{column}
  \end{columns}
\end{frame}

\subsection{Jeu de données}
\begin{frame}{\insertsubsection}
  \begin{columns}
    \begin{column}{.6\textwidth}
      \begin{itemize}
        \item Jeu de données issu de Open Food Facts
        \item[]
        \item Base de donnée ouverte et participative
        \item[]
        \item Repertorie les ingrédients, les allergènes, composition nutritionnelle, les labels, l'origine, etc
        \item[]
        \item Grosse base de données (environ \qty{4.5}{\giga\octet})
      \end{itemize}
    \end{column}
    \begin{column}{.4\textwidth}
      \centering
      \includegraphics[width=.8\textwidth]{./openfoodfacts-logo-en.pdf}
    \end{column}
  \end{columns}
\end{frame}

\section{Idée d'application}
\subsection{MIAM Iteractive Alimentation Manager}
\begin{frame}{\insertsubsection}
  \begin{columns}[t]
    \begin{column}{.4\textwidth}
      À partir des données présentent dans la base de données
      \begin{itemize}
        \item[]
        \item Pouvoir scanner un produit et trouver des alternatives similaires de meilleure qualité
              \begin{itemize}
                \item Meilleures pour la santé
                      \begin{itemize}
                        \item Avec moins d'additifs
                      \end{itemize}
                \item Meilleures pour l'environement
                      \begin{itemize}
                        \item Origine plus proche
                      \end{itemize}
                \item Voir les deux
              \end{itemize}
        \item[]
        \item Pouvoir choisir selon un régime alimentaire particulier
              \begin{itemize}
                \item Végétarien
                \item Vegan
                \item Halal
                \item Kasher
                \item Sans gluten
                \item etc
              \end{itemize}
      \end{itemize}
    \end{column}
    \begin{column}{.6\textwidth}
      Exemple du cas des saucisses :
      \begin{itemize}
        \item[]
        \item Nombre total d'aliments de ce type : 9120
      \end{itemize}
      \begin{table}
        \small
        \input{./Tableaux/SaucisseCat.tex}
      \end{table}
      \begin{table}
        \input{./Tableaux/SaucisseNbRegime.tex}
      \end{table}
    \end{column}
  \end{columns}
\end{frame}

\begin{frame}{\insertsubsection}
  Développer l'aspect participatif en incrémentant des données dans la base de données
  \begin{itemize}
    \item[]
    \item Beaucoup de lacune dans les données
          \begin{itemize}
            \item[] Implémenter un système de reconnaissance d'image
              permettant d'extraire les données manquantes des photos des étiquettes
          \end{itemize}
    \item[]
    \item Permettre de renseigner le prix de l'aliment et  de comparer le prix
          entre les aliments selon la qualité
          \begin{itemize}
            \item Création de listes par l'utilisateur pour évaluer le prix d'un panier
            \item[]
            \item Création de listes automatique d'aliments similaires de meilleure qualité à partir des listes établies
            \item[]
            \item Pouvoir comparer le prix des paniers/listes entre elles permettrait aux ménages de voir
                  la différence de budget entre différents paniers de qualités variées
                  et pouvoir choisir quel type d'alimentation il veulent/peuvent s'offrir
                  en fonction de la part du pouvoir d'achat qu'ils sont prêt à mettre dans l'alimentation
          \end{itemize}
  \end{itemize}
\end{frame}

\section{Nettoyage du jeu de données}
\subsection{Produits vendus en France}
\begin{frame}{\insertsubsection}
  \begin{columns}[T]
    \begin{column}{.4\textwidth}
      \begin{itemize}
        \item Sélection des produits vendus en France
        \item \num{850000} entrées sur les 2 milliards d'origine
      \end{itemize}
    \end{column}
    \begin{column}{.6\textwidth}
      \begin{itemize}
        \item Suppression des colonnes dans lesquelles on a moins de 1\% de données
              et des colonnes contenant des tags et des versions anglaises qui sont redondantes avec les
              colonnes du mêmes nom
        \item 57 colonnes restantes sur les 186 d'origine
      \end{itemize}
    \end{column}
  \end{columns}
  \vfill
  \begin{figure}
    \includegraphics[width=.75\textwidth]{DataDropFinal.pdf}
  \end{figure}
\end{frame}

\subsection{Nettoyage des outliers}
\begin{frame}{\insertsubsection}
  \begin{columns}
    \begin{column}{.3\textwidth}
      \begin{itemize}
        \item Valeurs extrêmes pour :
              \begin{itemize}
                \item Quantité servie
                \item Empreinte carbone
                \item Énergie
              \end{itemize}
        \item[]
        \item Des valeurs très importantes là où le nombre de valeurs est très faible
        \item[]
        \item Nous écartons ces valeurs lorsqu'elles sont en dehors du quartile qui contient 99\% des valeurs.
      \end{itemize}
    \end{column}
    \begin{column}{.7\textwidth}
      \begin{figure}
        \includegraphics[width=\textwidth]{subplotHistNettoyage.pdf}
      \end{figure}
    \end{column}
  \end{columns}
\end{frame}

\begin{frame}{\insertsubsection}
  \begin{columns}
    \begin{column}{.3\textwidth}
      \begin{itemize}
        \item On obtient des distributions plus intéressantes
        \item[]
        \item Fonctionnement assez similaire des différents types d'énergies
      \end{itemize}
    \end{column}
    \begin{column}{.7\textwidth}
      \begin{figure}
        \includegraphics[width=\textwidth]{subplotHistQ99.pdf}
      \end{figure}
    \end{column}
  \end{columns}
\end{frame}

\subsection{Données redondantes}
\begin{frame}{\insertsubsection}
  \begin{columns}
    \begin{column}{.4\textwidth}
      \begin{itemize}
        \item Mêmes valeurs energy\_100g et energy-kj\_100g
        \item[]
        \item Variation pour ce qui est de l'energie en kcal pour 100g
        \item[]
        \item Nous n'allons garder que les colonnes avec des unités d'énergie : kcal et kj
        \item[]
        \item De même nous allons supprimer les colonnes finissant par \_tags et \_en
              qui contiennent des données redondantes des colonnes du même nom sans suffixes
      \end{itemize}
    \end{column}
    \begin{column}{.6\textwidth}
      \begin{figure}
        \includegraphics[width=.9\textwidth]{EnergyReg.pdf}
      \end{figure}
    \end{column}
  \end{columns}
\end{frame}

\section{Exploration du jeu de données}
\subsection{Analyse univariée}
\begin{frame}{\insertsubsection}
  \begin{table}
    \tiny
    \input{./Tableaux/StatsNum1.tex}
    \input{./Tableaux/StatsNum2.tex}
    \input{./Tableaux/StatsNum3.tex}
    \caption{Tableaux des statistiques sur chaques variables numériques}
  \end{table}
\end{frame}

\begin{frame}{\insertsubsection}
  \begin{figure}
    \includegraphics[height=.8\textheight]{subplotBox}
  \end{figure}
\end{frame}

\begin{frame}{\insertsubsection}
  \begin{columns}
    \begin{column}{.5\textwidth}
      \begin{itemize}
        \item Observation du nombre d'occurences de chaques catégories de nutriscore
        \item[]
        \item Diagramme circulaire permet de visualiser les proportions
        \item[]
          \begin{itemize}
            \item \qty{50}{\percent} des aliments sont de catégorie D ou E
            \item \qty{75}{\percent} sont de catégorie C, D ou E
            \item seulemnt 1/4 des produits sont de catégorie A ou B
          \end{itemize}
        \item[]
        \item Il n'y a qu'une chance sur 4 de tirer un produit de bonne qualité au hasard
        \item[]
        \item Permettre au consommateur de mieux choisir ses aliments permettra d'améliorer la santé publique
      \end{itemize}
    \end{column}
    \begin{column}{.5\textwidth}
      \begin{figure}
        \includegraphics[width=\textwidth]{NutriPie.pdf}
      \end{figure}
    \end{column}
  \end{columns}
\end{frame}

\subsection{Analyse bivariée}
\begin{frame}{\insertsubsection}
  \begin{columns}
    \begin{column}{.55\textwidth}
      \begin{itemize}
        \item Les catégories du nutriscore semblent corrélées avec le score
        \item[]
        \item Vérification par ANOVA à une dimension
        \item H0 les moyennes des différentes catégories ne varient pas en fonction du score
      \end{itemize}
      \begin{table}
        \tiny
        \input{./Tableaux/ANOVANutriScoreGrade.tex}
        \caption{Résultats de l'ANOVA à une dimension entre le score du nutriscore et ses catégories}
      \end{table}
      \begin{itemize}
        \item p-value \qty{< 5}{\percent} on rejette l'hypothèse pour chaque catégories
      \end{itemize}
    \end{column}
    \begin{column}{.45\textwidth}
      \begin{figure}
        \includegraphics[width=\textwidth]{NutriBox.pdf}
      \end{figure}
    \end{column}
  \end{columns}
\end{frame}

\subsubsection{Matrice des corrélations}
\begin{frame}{\insertsubsubsection}
  \begin{columns}
    \begin{column}{.5\textwidth}
      \begin{itemize}
        \item La matrice des corrélations nous permet d'observer quelles variables
              semblent fonctionner de façon similaire
        \item[]
        \item Nous allons faire quelques regressions entre ces variables proches
      \end{itemize}
    \end{column}
    \begin{column}{.5\textwidth}
      \includegraphics[width=.98\textwidth]{HeatmapNum.pdf}
    \end{column}
  \end{columns}
\end{frame}

\subsubsection{Régressions}
\begin{frame}{\insertsubsubsection}
  \begin{columns}
    \begin{column}{.5\textwidth}
      \begin{figure}
        \includegraphics[width=\textwidth]{SelReg.pdf}
        \includegraphics[width=\textwidth]{EnergyFatReg.pdf}
      \end{figure}
    \end{column}
    \begin{column}{.5\textwidth}
      \begin{figure}
        \includegraphics[width=\textwidth]{SucreReg.pdf}
        \includegraphics[width=\textwidth]{GrasReg.pdf}
      \end{figure}
    \end{column}
  \end{columns}
\end{frame}

\subsection{Analyse multivariée : nutriscore et variables sucre, sel, gras et energie}
\begin{frame}{\insertsubsection}
  \begin{columns}
    \begin{column}{.5\textwidth}
      \begin{itemize}
        \item Régression multiple entre le nutriscore et les variables
              sucre, gras, sel et energie
        \item[]
        \item H0 les moyennes des variables sont égales entre elles
              ainsi qu'avec celle du nutriscore
        \item[]
        \item p-value \qty{< 5}{\percent} les variables ne sont donc pas très similaires
        \item[]
        \item Nous allons tout de même réaliser quelques graphiques des régressions
              entre Nutriscore et chaques variables
      \end{itemize}
    \end{column}
    \begin{column}{.5\textwidth}
      \begin{table}
        \tiny
        \input{./Tableaux/MultiReg}
      \end{table}
    \end{column}
  \end{columns}
\end{frame}

\subsubsection{Régressions avec le nutriscore}
\begin{frame}{\insertsubsubsection}
  \begin{columns}
    \begin{column}{.5\textwidth}
      \begin{figure}
        \includegraphics[width=\textwidth]{NutriEnergyReg.pdf}
        \includegraphics[width=\textwidth]{NutriSucreReg.pdf}
      \end{figure}
    \end{column}
    \begin{column}{.5\textwidth}
      \begin{figure}
        \includegraphics[width=\textwidth]{NutriGrasReg.pdf}
        \includegraphics[width=\textwidth]{NutriSelReg.pdf}
      \end{figure}
    \end{column}
  \end{columns}
\end{frame}

\subsubsection{ACP}
\begin{frame}{\insertsubsubsection}
  \begin{columns}
    \begin{column}{.5\textwidth}
      \begin{itemize}
        \item L'ACP permet de regrouper des variables qui fonctionnent de manière similaire
        \item[]
        \item ACP sur les colonnes numériques
        \item[]
        \item Le graphique des variances cumulées nous montre qu'il n'y a pas de rupture nette ("coude")
              dans la part expliquée des différentes composantes
      \end{itemize}
    \end{column}
    \begin{column}{.5\textwidth}
      \includegraphics[width=\textwidth]{ScreePlot.pdf}
    \end{column}
  \end{columns}
\end{frame}

\begin{frame}
  \begin{columns}
    \begin{column}{.5\textwidth}
      \includegraphics[width=\textwidth]{PCAF1F2.pdf}
    \end{column}
    \begin{column}{.5\textwidth}
      \includegraphics[width=\textwidth]{PCAF1F3.pdf}
    \end{column}
  \end{columns}
\end{frame}

\begin{frame}
  \begin{columns}
    \begin{column}{.5\textwidth}
      \includegraphics[width=\textwidth]{PCAF1F4.pdf}
    \end{column}
    \begin{column}{.5\textwidth}
      \includegraphics[width=\textwidth]{PCAF2F3.pdf}
    \end{column}
  \end{columns}
\end{frame}

\begin{frame}
  \begin{columns}
    \begin{column}{.5\textwidth}
      \includegraphics[width=\textwidth]{PCAF2F4.pdf}
    \end{column}
    \begin{column}{.5\textwidth}
      \includegraphics[width=\textwidth]{PCAF3F4.pdf}
    \end{column}
  \end{columns}
\end{frame}

\section{Conclusion}
\begin{frame}
  \begin{itemize}
    \item Jeu de données important
    \begin{itemize}
      \item Malgrès les lacunes il y a énormément d'informations sur les produits
      \item[]
      \item Permet de faire des choix précis au niveau du type d'alimentation dans l'application
    \end{itemize}
    \item[]
    \item Présences de liens vers des étiquettes dans la base de donnée permettrait d'entraîner un modèle de reconnaissance
    d'images afin d'implémenter la lecture des éléments manquants 
    \item[]
    \item L'ACP permet de mieux visualiser les liens entre variables
          \begin{itemize}
            \item Permet de catégoriser les aliments selon des critères plus larges que les variables
            \item[]
            \item Utile pour regrouper des produits et permettre de proposer des listes pertinentes dans l'application
          \end{itemize}
    \item[]
    \item Face aux maladies dites de civilisation (diabète, obésité, etc), permettre aux consommateurs
          de mieux se nourrir pourra participer à l'amélioration de la santé publique
  \end{itemize}
\end{frame}

\end{document}